\section{Parameter-free Algorithm}
\label{sec:DP-OTB_parameter-free}

In this section, we apply Algorithm~\ref{alg:DP-OTB:main} with a ``parameter-free'' online learner. These are algorithms that guarantee $\regret_T(u)\le \tilde O(\|u\|\sqrt{T})$ for all competitors $u$ simultaneously \citep{orabona2016coin, cutkosky2018black, mhammedi2020lipschitz}. By shifting coordinates, it is possible to obtain $\regret_T(u)\le \tilde O(\|u-x_0\|\sqrt{T})$ for any pre-specified point $x_0$. Thus, if $x_0\approx u$ is some good initialization, perhaps generated by pretraining, this bound yields significantly smaller regret than if we had used the worst-case diameter bound $\|u-x_0\|\le D$.

In order to obtain this refined bound with  privacy, we need to make a small modification to our conversion. For simplicity, we focus on Euclidean space with 2-norm, and we assume the distribution $\calD$ is in addition sub-Gaussian, i.e., for $R\sim \calD$,
\begin{align*}
    \P\{\sup_{\|a\|_2=1} \langle R, a\rangle \geq \epsilon \} \leq \exp\left(-\frac{\epsilon^2}{2\sigma^2}\right).
\end{align*}
In general, the proof extends to any Banach space and any distribution $\calD$ that concentrates on it.

In previous analysis (Equation \eqref{eq:DP-OTB:basic-decomposition-2}), we roughly bound $\|w_t-x^*\|\leq D$. However, in this section, we come up with a finer high probability bound that maintains a dependence on $\|w_t\|, \|x^*\|$. We then replace the loss $\ell_t(w)$ in Algorithm \ref{alg:DP-OTB:main} with a regularized loss $\ell_t(w)+\xi_t\|w\|_2+\nu_t\|w\|_2^2$, and we show that the new algorithm with regularized loss can achieve a parameter-free bound. The complete proof is presented in Appendix \ref{app:DP-OTB:parameter-free}.

\begin{theorem}
\label{thm:DP-OTB:parameter-free-2-norm-step-1}
Suppose w.r.t. 2-norm, $W$ is bounded by $D$ and $\ell$ is $G$-Lipschitz and $H$-smooth. Suppose $\calD$ is $(V,\alpha)$-RDP distribution and is $\sigma_\calD$-sub-Gaussian. If we set $\weight_t=t^3$ (i.e. $k=3$) and set $\sigma_t^2$ as defined in \eqref{eq:DP-OTB:RDP-noise}, then with probability at least $1-\delta$,
\begin{align*}
    f(x_T)-f(x^*) \leq \frac{4}{T^4}\left( \regret_T(x^*) + \sum_{t=1}^T \xi_t(\|w_t\|_2+\|x^*\|_2) + \nu_t(\|w_t\|_2^2+\|x^*\|_2^2) \right).
\end{align*}
where $C$ is a universal constant, $A=8\sqrt{2}C^2, A'=8\sqrt{d}\sigma_\calD C^2, \kappa=1+DH/G$, 
% $\Phi=\sqrt{\log(16dT\log(2\kappa T)/\delta)}$, 
and
\begin{align*}
    &\xi_t = AG\Phi t^{5/2} + A'(G+DH)\frac{\Phi\log_2(2T)t^2}{\rho} , \
    \nu_t = 28AH\Phi t^{5/2}, \ 
    \Phi = \sqrt{\log\frac{20dT\log(2\kappa T)}{\delta}}.
\end{align*}
\end{theorem}

\begin{theorem}
\label{thm:DP-OTB:parameter-free-main}
Following the assumptions and notations in Theorem \ref{thm:DP-OTB:parameter-free-2-norm-step-1}, if we replace $\ell_t(w)$ in Algorithm \ref{alg:DP-OTB:main} with regularized loss $\bar \ell_t(w) = \ell_t(w) + \xi_t\|w\|_2 + \nu_t\|w\|_2^2$ and denote the associated regret as $\overline \regret_t$, then with probability at least $1-\delta$, $\L(x_T)-\L(x^*)$ is bounded by:
\begin{align*}
    % \L(x_T)-\L(x^*) 
    % &\leq 
    \frac{4\overline \regret_T(x^*)}{T^4} + \frac{8A\|x^*\|(G+28\|x^*\|H)\Phi}{\sqrt{T}} + \frac{8A'\|x^*\|(G+DH)\Phi\log_2(2T)}{\rho T}.
\end{align*}
Moreover, with probability at least $1-\delta$, for all $t$ and for all $w\in W, \bar g_t \in \partial \bar \ell_t(w)$,
\begin{align*}
    \|\bar g_t\|_2 \leq Gt^3 + A(2G+57DH)\Phi t^{5/2} + 2A'(G+DH)\frac{\Phi\log_2(2T)t^2}{\rho}.
\end{align*}
\end{theorem}

\begin{remark}
If $\|\bar g_t\|_2 \leq \hat G$ for all $t$, parameter-free algorithms achieve regret bound $\regret_T(u) = \tilde O(\|u\|_2\hat G\sqrt{T})$. Therefore, with $\calD=\calN(0,I)$, $\calD$ is $1$-sub-Gaussian and $A'=O(\sqrt{d})$, so Theorem \ref{thm:DP-OTB:parameter-free-main} implies that with probability $1-\delta$,
\begin{align*}
    \frac{\overline\regret_T(x^*)}{T^4} = \tilde O\left( \frac{\|x^*\|_2G}{\sqrt{T}} + \frac{\|x^*\|_2(G+DH)}{T} + \frac{\|x^*\|_2\sqrt{d}(G+DH)}{\rho T^{3/2}} \right).
\end{align*}
Consequently,
\begin{align*}
    f(x_T) - f(x^*) \leq \tilde O\left( \frac{\|x^*\|_2(G+DH)}{\sqrt{T}} + \frac{\|x^*\|_2\sqrt{d}(G+DH)}{\rho T} \right).
\end{align*}
\end{remark}

\begin{proof}[Proof of Theorem \ref{thm:DP-OTB:parameter-free-main}]
The regularized regret satisfies
\begin{align*}
    \overline \regret_T(x^*) 
    % &= \sum_{t=1}^T [\ell_t(w_t) + \xi_t\|w_t\|_2 + \nu_t\|w_t\|_2^2] - [\ell_t(x^*)+\xi_t\|x^*\|_2^2 + \nu_t\|x^*\|_2^2] \\
    &= \regret_T(x^*) + \sum_{t=1}^T \xi_t(\|w_t\|_2-\|x^*\|_2)+\nu_t(\|w_t\|_2^2-\|x^*\|_2^2).
\end{align*}
Hence, upon substituting this equation into Theorem \ref{thm:DP-OTB:parameter-free-2-norm-step-1}, we get: with probability at least $1-\delta$,
\begin{align*}
    &f(x_T)-f(x^*) 
    \leq \frac{4\overline\regret_T(x^*)}{T^4} + \frac{8}{T^4} \sum_{t=1}^T \xi_t\|x^*\|_2 + \nu_t\|x^*\|_2^2 \\
    &\leq \frac{4\regret_T(x^*)}{T^4} + \frac{8(AG\Phi\|x^*\|_2+28AH\Phi\|x^*\|_2^2)}{\sqrt{T}} + \frac{8A'(G+DH)\Phi\log_2(2T)\|x^*\|}{\rho T}.
    % &\leq \frac{4\regret_T(x^*)}{T^4} + \frac{8A\|x^*\|(G+28\|x^*\|H)\Phi}{\sqrt{T}} + \frac{8A'\|x^*\|(G+DH)\Phi\log_2(2T)}{\rho T}.
\end{align*}
The second inequality is from $\sum_{t=1}^T \xi_t \leq T\xi_T$ because $\xi_t$ is increasing with $t$ (so is $\nu_t$). 

For the second part of the theorem, for each fixed $t$ and for all $\bar g_t\in \bar\ell_t(w_t)$, $\bar g_t = g_t+\gamma_t + \xi_t u+ 2\nu_t w_t$, where $u\in \partial \|w_t\|_2$ and thus $\|u\|_2 \leq 1$. Therefore,
\begin{align*}
    \|\bar g_t\|_2
    &\leq \| g_t - \weight_t\nabla f(x_t)\|_2 + \|\weight_t\nabla f(x_t)\|_2 + \|\gamma_t\|_2 + \|\xi_t u\|_2 + \|2\nu_tw \|_2
\end{align*}
By Lipschitzness, $\|\weight_t\nabla f(x_t)\|_2 \leq Gt^3$. Since $W$ is bounded, $\|2\nu_tw_t\|_2\leq 2D\nu_t$. Also, $\|\xi_tu\|_2\leq \xi_t$. Moreover, we can prove (in Eq. \ref{eq:DP-OTB:variance-concentration} and \ref{eq:DP-OTB:privacy-concentration}) that for each $t$, with probability at least $1-\delta/2T$,
\begin{align*}
    & \|\weight_t\nabla f(x_t) - g_t\|_2
    \leq 8C^2\Phi \sqrt{\sum_{i=1}^t i^4(G+H\|w_i-x_{i-1}\|_2)^2} 
    \leq A\Phi (G+DH) t^{5/2}, \\
    & \|\gamma_t\|_2 \leq \frac{A'}{\rho}(G+DH)\Phi \log_2(2T)t^2.
\end{align*}
Upon taking the union bound for all $t$ and  the definition of $\xi_t,\nu_t$, we get the desired bound.
% For the third term, we can prove (in Equation \eqref{eq:DP-OTB:privacy-concentration}) that for each $t$, with probability at least $1-\delta/2T$,
% \begin{align*}
%     \|\gamma_t\|_2 \leq \frac{A'}{\rho}(G+DH)\Phi \log_2(2T)t^2.
% \end{align*}
% In conclusion, upon taking the union bound over all $t$ and substituting the definition of $\xi_t,\nu_t$, we get: with probability at least $1-\delta$, for all $t$,
% \begin{align*}
%     \| \bar g_t \|_2 
%     &\leq A\Phi(G+DH)t^{5/2} + Gt^3 + A'\Phi(G+DH) \frac{\log_2(2T)t^2}{\rho} \\
%     &\quad \ + \left(A\Phi Gt^{5/2} +  A'\Phi(G+DH)\frac{\log_2(2T)t^2}{\rho}\right) + 56A\Phi DHt^{5/2} \\
%     &\leq Gt^3 + A(2G+57DH)\Phi t^{5/2} + 2A'(G+DH)\frac{\Phi\log_2(2T)t^2}{\rho}.
% \end{align*}
\end{proof}